% Additive standalone worked exercise for original ch017.
\clearpage
\section*{Verification case: an array is a measurement with a contract}
\noindent A mathematical vector has a dimension; a software array additionally has
a shape, storage type, stride, coordinate convention and validity policy.
For a point cloud \(P\in\mathbb R^{N\times 3}\), broadcasting a translation
\(t\in\mathbb R^3\) across its rows implements \(P'=P+\mathbf1_Nt^T\).
This is correct only if each row is a point in the \emph{same}
frame and unit system as \(t\), not simply because the array shapes permit it.

For an inspectable experiment, choose the three points (in metres)
\[
P=\begin{pmatrix}1&0&0\\0&2&0\\0&0&3\end{pmatrix},
\qquad t=\begin{pmatrix}0.5\\-1\\2\end{pmatrix}.
\]
The translated rows are respectively
\((1.5,-1,2)\), \((0.5,1,2)\), and \((0.5,-1,5)\).
An equivalent C++23 implementation may loop over rows, while NumPy
may broadcast \(t\); the expected geometry does not depend on the
language or execution hardware.

\paragraph{Limits of the test.}
The three expected points verify the chosen examples, not arbitrary
input. A correct implementation also rejects the wrong number of
components, incompatible layout assumptions and nonfinite values
when the application's contract requires finite geometry.
Integer arrays can silently truncate a fractional translation after an
unsafe cast. Float32 and Float64 may differ because of rounding.
Whether a view avoids a copy is a question about memory and ownership,
not a theorem about broadcasting.

\paragraph{Reconstruction.}
Implement the update in NumPy and in a documented C++23 loop;
compare after converting both to a declared precision. Include
three negative tests (shape, dtype and nonfinite input), a
translation-inverse round trip, and a test confirming that point
differences do not change under a uniform translation. Record any
failure to preserve the latter within a stated floating-point
tolerance rather than silently treating a visual overlap as proof.
